\clearpage\section*{EV2-B -- Blanc 1}\label{EV2-B}
\addcontentsline{toc}{section}{EV2-B -- Blanc 1}
\textbf{Sujet réservé au professeur}. Chapitres 1 à 7.\par Durée : 2 h. Sans calculatrice. Note sur 20. 14–16 janvier 2027.
\subsection*{Partie 1 -- Automatismes (6 points)}
Douze questions, 0,5 point chacune. Une seule réponse correcte ; aucune justification demandée. Une réponse incorrecte ou absente rapporte 0, sans retrait de points.
\begin{enumerate}
\item Calculer $5+2\times3$.\par A. $11$\quad B. $10$\quad C. $14$\quad D. $21$
\item Calculer $\frac12+\frac13$.\par A. $\frac{5}{3}$\quad B. $\frac{1}{6}$\quad C. $\frac{5}{6}$\quad D. $\frac{2}{5}$
\item Calculer $\dfrac{2^5}{2^{3}}$.\par A. $2$\quad B. $10$\quad C. $1$\quad D. $4$
\item Une valeur de 50 augmente de 20 \%. Valeur finale ?\par A. $10$\quad B. $60$\quad C. $40$\quad D. $70$
\item Une valeur de 80 est obtenue après une baisse de 20 \%. Valeur initiale ?\par A. $100$\quad B. $64$\quad C. $96$\quad D. $60$
\item Hausse de 25 \% puis baisse de 20 \% : coefficient global ?\par A. $\frac{9}{10}$\quad B. $1$\quad C. $\frac{21}{20}$\quad D. $\frac{4}{5}$
\item La solution de $2x+5=0$ vaut\par A. $-10$\quad B. $\frac{5}{2}$\quad C. $\frac{-5}{2}$\quad D. $10$
\item La pente entre $(0;5)$ et $(2;11)$ vaut\par A. $2$\quad B. $-3$\quad C. $3$\quad D. $6$
\item La moyenne de $5, 7, 9$ vaut\par A. $7$\quad B. $21$\quad C. $8$\quad D. $5$
\item Un secteur de $90^\circ$ représente quelle fraction du total ?\par A. $\frac{1}{3}$\quad B. $\frac{1}{2}$\quad C. $\frac{3}{4}$\quad D. $\frac{1}{4}$
\item Une probabilité vaut $\dfrac{2}{5}$. Son complément vaut\par A. $\frac{5}{2}$\quad B. $\frac{2}{5}$\quad C. $\frac{1}{5}$\quad D. $\frac{3}{5}$
\item Pour $f(x)=x^2+5$, $f(-2)$ vaut\par A. $1$\quad B. $7$\quad C. $9$\quad D. $20$
\end{enumerate}
\clearpage\subsection*{Partie 2 -- Exercices indépendants (14 points)}
Toute démarche doit être justifiée. Les résultats exacts sont attendus sauf indication contraire.
\Needspace{12\baselineskip}
\subsection*{Exercice 1 -- Indépendance et répétition (5 points)}
Deux événements $A$ et $B$ sont indépendants et vérifient $P(A)=\dfrac14$ et $P(B)=\dfrac23$. La première question porte sur ces événements. Dans une expérience distincte, utilisée aux questions 2 et 3, on répète trois fois de manière indépendante une épreuve dont la probabilité de succès vaut $\dfrac14$.
\begin{enumerate}
\item Calculer $P(A\cap B)$ et $P(A\cup B)$. \hfill(1 pt)
\item Calculer la probabilité d’exactement deux succès dans les trois essais. \hfill(2 pts)
\item Calculer celle d’au moins un succès et justifier le passage au complémentaire. \hfill(2 pts)
\end{enumerate}
\Needspace{12\baselineskip}
\subsection*{Exercice 2 -- Rationnel et symétrie (5 points)}
La fonction $f$ est définie sur $\R\setminus\{0\}$ par $f(x)=\dfrac{x^2+4}{x}$. Son domaine est symétrique par rapport à 0, mais ce n’est pas un intervalle. On s’appuiera sur sa parité pour relier les études sur les deux demi-droites.
\begin{enumerate}
\item Étudier la parité et réécrire f. \hfill(1 pt)
\item Étudier les variations sur $]0;+\infty[$. \hfill(2 pts)
\item En déduire les variations sur $]-\infty;0[$ et expliquer pourquoi f n’est pas globalement croissante sur son domaine. \hfill(2 pts)
\end{enumerate}
\Needspace{12\baselineskip}
\subsection*{Exercice 3 -- Distance et paramètre (4 points)}
Dans un repère orthonormé, on donne la droite $d$ d’équation $3x+4y-12=0$ et le point $P(0;0)$. On note $H$ le projeté orthogonal de $P$ sur $d$. On cherche le cercle de centre $P$ tangent à cette droite.
\begin{enumerate}
\item Donner une normale et vérifier que $(4;0)$ est sur d. \hfill(1 pt)
\item Calculer le projeté H de P. \hfill(1 pt)
\item Déterminer la distance de P à d et l’équation du cercle de centre P tangent à d. \hfill(2 pts)
\end{enumerate}
